\subsection{自由群的下中心列}

我们将证明下面两个定理：

定理 2. 假设在环 K 中，关系 n.1 = 0 对每个整数 n 蕴涵 n = 0。设 r 是 X 到 \(\hat{\mathrm{A}}(\mathrm{X})\) 中的一个映射，使得 \(\omega(r(x)) \geqslant 2\) （对 \(x \in X\) ），并设 g 是 \(\mathrm{F}(\mathrm{X})\) 到 Magnus 群 \(\Gamma(\mathrm{X})\) 中的一个同态，使得 \(g(x) = 1 + x + r(x)\) （对 \(x \in X\) ）。对一切 \(n \geqslant 1\) ， \(\mathrm{C}^n\mathrm{F}(\mathrm{X})\) 是 g 之下子群 \(1 + \hat{\mathrm{A}}_n(\mathrm{X})\) （\(\Gamma(\mathrm{X})\) 的）的逆像。

定理 3. 对一切 \(x \in \mathbf{X}\) ，设 \(c(x)\) 是 \(x\) 在 \(\mathrm{F}(\mathbf{X}) / (\mathrm{F}(\mathbf{X}), \mathrm{F}(\mathbf{X}))\) 中的典范像。设 \(\mathfrak{g}\) 是分次李 \(\mathbf{Z}\) -代数，它与滤过 \((\mathrm{C}^{\mathrm{n}}\mathrm{F}(\mathbf{X}))_{n \geqslant 1}\) （\(\mathrm{F}(\mathbf{X})\) 的）相伴 (§ 4， no. 6)。自由李 \(\mathbf{Z}\) -代数 \(\mathrm{L}_{\mathbf{Z}}(\mathbf{X})\) 到 \(\mathfrak{g}\) 中的唯一同态——它扩张 \(c\) ——是一个同构。

粗略地说，与自由群 \(F(X)\) （带下中心列）相伴的分次李 Z-代数就是自由李 Z-代数 \(L_{Z}(X)\) 。

我们记 \(\mathrm{F}(\mathbf{X}) = \mathrm{F}\) 、 \(\Gamma(\mathbf{X}) = \Gamma\) 、 \(\hat{\mathrm{A}}(\mathbf{X}) = \hat{\mathrm{A}}\) 、 \(\hat{\mathrm{A}}_{\mathbf{Z}}(\mathbf{X}) = \hat{\mathrm{A}}_{\mathbf{Z}}\) 、 \(\mathrm{C}^n\mathrm{F}(\mathbf{X}) = \mathrm{C}^n\) 、 \(\Gamma_n = 1 + \hat{\mathrm{A}}_n(\mathbf{X})\) ，并设 \(\alpha: \mathrm{L}_{\mathbf{Z}}(\mathbf{X}) \to \mathfrak{g}\) 是定理 3 的陈述中引入的同态。

\textbf{(A) 预备性化简。}

用 \(\gamma\) 表示 F 到 \(\Gamma\) 中的一个同态，它由 \(\gamma(x) = 1 + x\) （对 \(x \in \mathbf{X}\) ）所定义。由引理 4，存在自同构 \(\sigma\) （代数 \(\hat{\mathbf{A}}\) 的），它与 \(\hat{\mathbf{A}}\) 上的滤过相容，并且使得 \(\sigma(1 + x) = g(x)\) （对一切 \(x \in \mathbf{X}\) ）；于是 \(\sigma(\Gamma_n) = \Gamma_n\) （对一切 \(n\) ）。由于同态 \(g\) 与 \(\sigma \circ \gamma\) （都是从 F 到 \(\Gamma\) 中）在 \(\mathbf{X}\) 上一致，故 \(g = \sigma \circ \gamma\) ，从而 \(g^{-1}(\Gamma_n) = \gamma^{-1}(\Gamma_n)\) 。在定理 2 的假设下， \(\mathbf{Z}\) 可以与 \(\mathbf{K}\) 的一个子环等同；因此 Magnus 代数 \(\hat{\mathbf{A}}_{\mathbf{Z}}\) 与 \(\hat{\mathbf{A}}\) 的一个子环等同，并且 \(\hat{\mathbf{A}}_{\mathbf{Z}}\) 上的滤过由 \(\hat{\mathbf{A}}\) 上的滤过诱导。由于 \(\gamma\) 把 F 映到 \(\hat{\mathbf{A}}_{\mathbf{Z}}\) 中，我们看出只需在附加假设 \(\mathbf{K} = \mathbf{Z}\) 、 \(r = 0\) 从而 \(g = \gamma\) 之下证明定理 2 与定理 3，我们今后都作这些假设。

\textbf{(B) \(\alpha\) 的满射性。}

由于 \(\mathbf{X}\) 生成群 \(\mathbf{F} = \mathbf{C}^1\) ，集合 \(c(\mathbf{X})\) 生成 \(\mathbf{Z}\) -模 \(\mathfrak{g}^1 = \mathbf{C}^1 / \mathbf{C}^2\) 。而 \(\mathfrak{g}^1\) 生成李 \(\mathbf{Z}\) -代数 \(\mathfrak{g}\) (§ 4， no. 6， 命题 5)，因而 \(c(\mathbf{X})\) 生成 \(\mathfrak{g}\) ，这就证明了 \(\alpha\) 是满射。

\begin{enumerate}
\def\labelenumi{(\Alph{enumi})}
\setcounter{enumi}{2}
\tightlist
\item
  我们把分次代数 \(\mathrm{gr}(\hat{\mathbf{A}})\) 与 A(X) 等同，所用的典范同构为 \(\mathrm{A}^n (\mathbf{X})\to \hat{\mathrm{A}}_n / \hat{\mathrm{A}}_{n + 1}\) 。对每个整数 \(n\geqslant 1\) ，我们记 \(\mathrm{F^n} = \gamma^{-1}(\Gamma_n)\) ；我们知道 (§ 4， no. 5) \((\mathbf{F}^n)_{n\geqslant 1}\) 是 F 上的一个整中心滤过。用 \(\mathfrak{g}'\) 表示相伴的分次李 Z-代数 (§ 4， no. 4)。设 \(f\) 是李代数同态，它从 \(\mathfrak{g}'\) 到 A(X) 中，与 \(\gamma\) 相伴 (§ 4， no. 5， 命题 3)。而 \(\mathrm{C^n}\subset \mathrm{F^n}\) （对每个整数 \(n\geqslant 1\) ）(§ 4， no. 6， 命题 4)，因而存在一个典范同态 \(\varepsilon\) ，从 \(\mathfrak{g} = \bigoplus_{n\geqslant 1}\mathrm{C}^n /\mathrm{C}^{n + 1}\) 到 \(\mathfrak{g}' = \bigoplus_{n\geqslant 1}\mathrm{F^n / F^{n + 1}}\) 中
\end{enumerate}

\[
\mathrm{L} _ {\mathbf {Z}} (\mathrm{X}) \stackrel {{\alpha}} {{\longrightarrow}} \mathfrak {g} \stackrel {{\varepsilon}} {{\longrightarrow}} \mathfrak {g} ^ {\prime} \stackrel {{f}} {{\longrightarrow}} \mathrm{A} (\mathrm{X}).
\]

我们记 \(\beta = f\circ \varepsilon\) ；我们如下明确地给出 \(\beta\) ：若 \(u\) 是模 \(C^{n + 1}\) 的类（元素 \(w\) 取自 \(C^n\) ），则 \(\gamma (w) - 1\) 有阶 \(\geqslant n\) （在 \(\hat{\mathbf{A}}\) 中），并且 \(\beta (u)\) 是 \(\gamma (w) - 1\) 的 \(n\) 次齐次分量。特别地，

\[
\beta (c (x)) = x \quad \text { for   all } x \in \mathbf {X}.\tag{3}
\]
% semantic-fragments: legacy-input-seam
\relax{}
\textbf{(D) 定理 2 和 3 的证明。}

李代数同态 \(\beta \circ \alpha: \mathrm{L}_{\mathbf{Z}}(\mathrm{X}) \to \mathrm{A}(\mathrm{X})\) 限制于 \(\mathbf{X}\) 后，由 (3) 可知是恒等映射，因而是典范嵌入（§ 3，第 1 号）。所以 \(\alpha\) 是单射，并且由 (B) 是双射；这证明了定理 3。由于 \(\beta \circ \alpha = f \circ \varepsilon \circ \alpha\) 是单射且 \(\alpha\) 是双射，\(\varepsilon\) 是单射。对所有 \(n \geqslant 1\)，

\[
\varepsilon_ {n} \colon \mathrm{C} ^ {n} / \mathrm{C} ^ {n + 1} \rightarrow \mathrm{F} ^ {n} / \mathrm{F} ^ {n + 1}
\]

是单射，因而

\[
\mathrm{C} ^ {n} \cap \mathrm{F} ^ {n + 1} = \mathrm{C} ^ {n + 1}.
\]

\(C^{1} = F = F^{1}\)；若 \(C^{n} = F^{n}\)，则 \(C^{n} \cap F^{n+1} = F^{n+1}\)，从而 \(C^{n+1} = F^{n+1}\)，这通过对 \(n \geqslant 1\) 作归纳证明了定理 2。

推论。

\[
\bigcap_ {n \geqslant 1} \mathrm{C} ^ {n} \mathrm{F} (\mathrm{X}) = \{e \}.
\]

取 K = Z 且 r = 0，应用定理 2，

\[
\bigcap_ {n \geqslant 1} \mathrm{C} ^ {n} \mathrm{F} (\mathrm{X}) = \bigcap_ {n \geqslant 1} \bar {g} (1 + \hat {\mathrm{A}} _ {n} (\mathrm{X})) = \bar {g} \left(\bigcap_ {n \geqslant 1} (1 + \hat {\mathrm{A}} _ {n} (\mathrm{X}))\right) = \bar {g} (1) = \{e \}.
\]

注。设 H 是关于 X 的 Hall 集（§2，第 10 号）。设 M 是在 F(X) 上由复合律 \((x,y)\mapsto (x,y) = x^{-1}y^{-1}xy\) 定义的二元代数系统，并设 \(\phi\) 是从 M(X) 到 M 的同态，其在 X 上的限制为恒等映射。称 \(\phi (\mathrm{H})\) 的元素为与 Hall 集 H 相关的 F(X) 的基本换位子。对每个整数 \(n\geqslant 1\)，令 \(\mathrm{H}_{\mathfrak{n}}\) 为 H 的子集，它由长度为 \(n\) 的元素组成；我们知道（§2，第 11 号，定理 1），\(\mathrm{H}_{\mathfrak{n}}\) 到 \(\mathbf{L}_{\mathbf{Z}}(\mathbf{X})\) 的典范映射是阿贝尔群 \(\mathbf{L}_{\mathbf{Z}}^{n} (\mathbf{X})\) 的一组基。此外，\(\phi (\mathrm{H_n})\subset \mathrm{C^n}\)；对于所有 \(m\in \mathrm{H}_{\mathfrak{n}}\)，令 \(\phi_{\mathfrak{n}}(m)\) 表示 \(\mathrm{C}^{n + 1}\) 模 \(\phi (m)\in \mathrm{C}^n\) 的陪集。于是定理 3 表明，\(\phi_{\mathfrak{n}}\) 是从 \(\mathrm{H}_{\mathfrak{n}}\) 到阿贝尔群 \(\mathrm{C}^n /\mathrm{C}^{n + 1}\) 的一组基的双射。立即得到，对于所有 \(w\in F(\mathbf{X})\) 和所有 \(i\geqslant 1\)，存在 \(\alpha_{i}\) 中唯一的元素 \(\mathbf{Z}^{(\mathrm{H}_{i})}\)，使得对于 \(n\geqslant 1\)，

\[
w = \prod_ {i = 1} ^ {n} \prod_ {m \in H _ {i}} \phi (m) ^ {\alpha_ {i} (m)} \mod C ^ {n + 1},\tag{4}
\]

其中，乘积按照 H 上给出的全序计算。

例。设 X 是含有两个元素 x、y 的集合，并令 \(H_{1} = \{x, y\}\)、\(H_{2} = \{xy\}\)。因而 F(X) 的每个元素 w 都可以写成

\[
w \equiv x ^ {a} y ^ {b} (x, y) ^ {c} \bmod . \mathrm{C} ^ {3} \quad \text { with } a, b, c \text { in } \mathbf {Z}.
\]

对于 \(w = (xy)^n\)，有 \(a = b = n\) 且 \(c = n(1 - n)/2\)（参见习题 9），从而

\[
(x y) ^ {n} \equiv x ^ {n} y ^ {n} (x, y) ^ {n (1 - n) / 2} \mod C ^ {3}.
\]

